% =============================================================================
% Aula de Revisão — Unidade II Completa (04/05/2026)
% Semana 10 | Aula concentrada na segunda-feira (06/05 é feriado de Vitória de Santo Antão)
% Objetivo: cobrir semanas 08–12 em 3h + preparar alunos para Prova II, Final e II Chamada
% =============================================================================

\chapter{Revisão Integrada — Unidade II (04/05/2026)}

% ---------------------------------------------------------------------------
\section{Objetivo e Contexto}
% ---------------------------------------------------------------------------

\begin{NoteBox}
\textbf{Contexto desta aula:} as semanas 09 a 11 não foram ministradas presencialmente.
Esta aula de 3 horas concentra \textbf{todo o conteúdo da Unidade II} que cai nas provas
de II Unidade, Final e II Chamada — com toda a teoria essencial e questões no exato
formato cobrado nas avaliações.

\textbf{Calendário restante:}
\begin{itemize}
  \item \textbf{04/05} (hoje) — Revisão Integrada Unidade II (esta aula)
  \item \textbf{11/05} — Evento institucional (sem aula)
  \item \textbf{18/05} — Apresentação em grupo: Dashboard Storytelling (última aula)
  \item \textbf{25/05} — Início da semana de Provas (Prova II Unidade)
\end{itemize}
\end{NoteBox}

\begin{BoardBox}
\textbf{Roteiro da aula (19h–22h)}
\begin{enumerate}
  \item 19h–19h35 — \textbf{Bloco I:} Big Idea, AED vs Explanatória, Storyboard (S08)
  \item 19h35–20h05 — \textbf{Bloco II:} Design de Informação — Pré-atenção, Chartjunk, Relance (S09)
  \item 20h05–20h35 — \textbf{Bloco III:} Tableau Interativo — Filtros, Parâmetros e Ações (S10)
  \item 20h35–21h05 — \textbf{Bloco IV:} Dashboards Decisórios, KPIs e Pitch de 3 min (S11)
  \item 21h05–21h50 — \textbf{Simulado:} 5 questões no formato das provas (resolvidas em sala)
  \item 21h50–22h — \textbf{Orientações:} grupos e entrega da Semana 18
\end{enumerate}
\end{BoardBox}

% =============================================================================
\section{Bloco I — Big Idea, AED vs Explanatória e Storyboard (S08)}
% =============================================================================

\begin{SolvedBox}
\textbf{A virada de perspectiva: AED vs Análise Explanatória}

\begin{tabular}{|l|p{5.5cm}|p{5.5cm}|}
\hline
 & \textbf{AED (Exploratória)} & \textbf{Explanatória} \\
\hline
\textbf{Quem usa} & O analista & O tomador de decisão \\
\textbf{Objetivo} & Entender o dado, achar padrões & Provocar uma ação específica \\
\textbf{Linguagem} & ``O dado mostra X'' & ``O que este dado muda no que você faz?'' \\
\textbf{Visuais} & Centenas (a maioria descartada) & Apenas os que servem à mensagem \\
\textbf{Erro clássico} & — & Apresentar todo o processo exploratório para o CEO \\
\hline
\end{tabular}
\end{SolvedBox}

\begin{BoardBox}
\textbf{Big Idea Framework — 3 componentes obrigatórios (para o quadro)}

\begin{enumerate}
  \item \textbf{Ponto de vista:} afirmação, não observação neutra.
  (``O dado mostra X'' $\neq$ Big Idea. ``Devemos fazer Y porque X'' = Big Idea.)
  \item \textbf{Tensão:} o que está em jogo? Qual o risco de não agir?
  \item \textbf{Stake:} por que a audiência deveria se importar? O que a afeta diretamente?
\end{enumerate}

\textbf{Fórmula rápida:} [Recomendação de ação] + [dado concreto que justifica] + [impacto financeiro/operacional se nada for feito].
\end{BoardBox}

\begin{SolvedBox}
\textbf{Big Ideas boas vs. ruins — exemplos rápidos:}

\begin{tabular}{|p{5.5cm}|p{6.5cm}|}
\hline
\textbf{Ruim (observação neutra)} & \textbf{Boa (Big Idea com tensão)} \\
\hline
``A taxa de cancelamento subiu 12\%.'' &
``Se não corrigirmos o onboarding até setembro, perderemos R\$1,2M em receita recorrente.'' \\
\hline
``Canal online representa 60\% das vendas.'' &
``O canal físico gera 40\% das vendas com 70\% do custo — precisamos decidir sua continuidade antes do Q3.'' \\
\hline
``O churn é de 18\% ao ano.'' &
``Implementando contato proativo a cada 60 dias, reduzimos o churn de 18\% para 8\% e economizamos R\$185k/ano.'' \\
\hline
\end{tabular}

\textbf{Teste da Big Idea:} leia para alguém que não viu os dados. Se a resposta
for ``interessante'' sem que a pessoa saiba o que fazer — a Big Idea não está pronta.
\end{SolvedBox}

\begin{BoardBox}
\textbf{Storyboard — estrutura narrativa para o quadro}

Toda boa apresentação de dados segue 4 atos (Nancy Duarte):
\begin{enumerate}
  \item \textbf{Contexto} — estado atual com dados reais (KPIs)
  \item \textbf{Conflito} — a tensão: o que está errado ou em risco
  \item \textbf{Resolução} — sua proposta baseada em evidências
  \item \textbf{Call to Action} — decisão específica pedida, com prazo
\end{enumerate}

\textbf{Traduzindo para telas:}
Tela 1 (KPIs atuais) → Tela 2 (problema/oportunidade) → Tela 3 (análise) → Tela 4 (recomendação + próximo passo).
\end{BoardBox}

% ---------------------------------------------------------------------------
\subsection{Questão Tipo Prova I — Big Idea (S08)}
% ---------------------------------------------------------------------------

\begin{SolvedBox}
\textbf{Questão Q-A1 — Storytelling com Dados e Big Idea}

Felipe é analista de dados em uma rede de 12 restaurantes em Fortaleza. Após analisar
os dados operacionais, ele descobriu que:
\begin{itemize}
  \item 65\% das reclamações ocorrem no horário do almoço (12h–14h)
  \item Funcionários com menos de 60 dias cometem 3$\times$ mais erros do que os veteranos
  \item Unidades com supervisor fixo têm NPS 48\% maior que as sem supervisor
\end{itemize}

Qual das alternativas é a \textbf{Big Idea bem formulada} que Felipe deveria usar na
apresentação para a diretoria?

\begin{enumerate}[label=\Alph*)]
  \item ``Devemos implantar supervisão fixa e tutoria nos primeiros 60 dias de todo
  funcionário contratado, pois isso pode reduzir reclamações no almoço em até 55\%
  e elevar o NPS médio de 42 para 62 pontos.''
  \item ``Os dados mostram informações sobre reclamações e satisfação dos clientes.''
  \item ``Analisamos os registros de 12 restaurantes no período de 2024–2025.''
  \item ``O NPS médio da rede é de 42 pontos.''
  \item ``O gráfico mostra a distribuição de reclamações por horário.''
\end{enumerate}

\textbf{Gabarito: A.}

\textbf{Justificativa:} a alternativa A combina os 3 elementos obrigatórios da Big Idea:
(1) ponto de vista (``devemos implantar''), (2) dado concreto que justifica (``65\% das
reclamações, funcionários novos 3$\times$ piores'') e (3) impacto mensurável (``NPS de 42
para 62''). As demais são vagas~(B), meramente descritivas~(C, D) ou apenas rótulo de
gráfico~(E).
\end{SolvedBox}

% =============================================================================
\section{Bloco II — Design de Informação: Pré-atenção, Chartjunk e Teste do Relance (S09)}
% =============================================================================

\begin{SolvedBox}
\textbf{Atributos pré-atencionais — o cérebro lê antes de você pensar}

O cérebro humano detecta certos atributos visuais em menos de 250ms, sem esforço consciente.
Use esses atributos para guiar o olhar até o dado importante:

\begin{tabular}{|l|p{4.5cm}|p{5.5cm}|}
\hline
\textbf{Atributo} & \textbf{O que faz} & \textbf{Regra de uso} \\
\hline
\textbf{Cor} & Agrupa, destaca, emociona & 1 cor de destaque; resto em cinza \\
\textbf{Tamanho} & Indica magnitude & Maior = mais importante (consistente!) \\
\textbf{Posição} & Organiza hierarquia & Mais importante no canto sup. esquerdo \\
\textbf{Intensidade} & Contraste claro-escuro & Mais escuro = mais importante \\
\textbf{Enclosure} & Agrupa com borda/área & Caixa ao redor de elementos relacionados \\
\hline
\end{tabular}

\textbf{Regra de ouro:} use no máximo \textbf{2 atributos pré-atencionais} por visual.
Acima disso, a hierarquia colapsa — o olho não sabe para onde ir.
\end{SolvedBox}

\begin{BoardBox}
\textbf{Teste do Relance (3 segundos) — para o quadro}

Mostre o gráfico por 3 segundos, cubra. Pergunte:
\begin{enumerate}
  \item O que o gráfico mostra?
  \item Qual é a mensagem principal?
  \item O que você faria com essa informação?
\end{enumerate}

Se a audiência não responder à Q2 em 3 segundos, o visual falhou — não o leitor.

\textbf{Causas de falha mais comuns:}
\begin{itemize}
  \item Título rótulo (``Dados de Vendas 2025'') — solução: título mensagem
  \item Excesso de elementos — solução: remover tudo que não é dado
  \item Cores sem propósito — solução: 1 cor de destaque, resto cinza
  \item Legenda separada — solução: rótulo direto no dado
\end{itemize}
\end{BoardBox}

\begin{SolvedBox}
\textbf{Títulos rótulo vs títulos mensagem:}

\begin{tabular}{|l|l|}
\hline
\textbf{Título rótulo (ruim)} & \textbf{Título mensagem (bom)} \\
\hline
``Vendas por Região'' & ``Região Sul concentra 58\% das vendas'' \\
``Evolução Mensal'' & ``Churn caiu 34\% após o programa de fidelidade'' \\
``Satisfação de Clientes'' & ``NPS Premium é 2$\times$ superior ao do plano Básico'' \\
``Comparativo de Custos'' & ``Canal físico custa R\$12 por venda; digital, R\$0,80'' \\
\hline
\end{tabular}

\textbf{Fórmula:} [Sujeito] + [verbo de tendência/comparação] + [magnitude/contexto].
\end{SolvedBox}

\begin{BoardBox}
\textbf{Data-ink ratio (Tufte) — para o quadro}

\[
\text{Data-ink ratio} = \frac{\text{Tinta usada para representar dados}}{\text{Tinta total do gráfico}}
\]

Maximize esse ratio. \textbf{Chartjunk} = tudo que ocupa tinta sem carregar informação.
Remova: grades pesadas, barras 3D, sombras, gradientes sem propósito, legendas redundantes.

\textbf{Princípio:} ``Se eu remover isso, o gráfico fica menos claro?'' Se não, remova.
\end{BoardBox}

% ---------------------------------------------------------------------------
\subsection{Questão Tipo Prova II — Design de Dashboard (S09)}
% ---------------------------------------------------------------------------

\begin{SolvedBox}
\textbf{Questão Q-A2 — Design de Dashboards e Hierarquia Visual}

Camila é analista de BI em uma transportadora com 80 filiais no Nordeste. O diretor
de operações tem apenas 5 minutos por dia para analisar o dashboard. Ela está decidindo
entre duas abordagens de design:

\textbf{Opção A:}
\begin{itemize}
  \item Título: ``Atrasos no Q1 2026: 22\% acima da meta — foco na rota Fortaleza–Recife''
  \item KPI central grande mostrando ``122\%'' em vermelho (acima da meta de atrasos)
  \item Gráfico de barras: rotas com a rota crítica destacada em azul-escuro
  \item Cores: fundo branco, 2–3 cores com propósito; espaço em branco generoso
\end{itemize}

\textbf{Opção B:}
\begin{itemize}
  \item Título: ``Dashboard de Logística''
  \item 14 gráficos diferentes (pizza, barras, linhas, área, mapa, etc.)
  \item Todas as 9 cores da paleta padrão utilizadas simultaneamente
  \item Sem espaço em branco — ``aproveitando'' toda a tela
  \item Títulos genéricos: ``Gráfico 1'', ``Gráfico 2''…
\end{itemize}

Qual alternativa avalia \textbf{corretamente} as duas opções?

\begin{enumerate}[label=\Alph*)]
  \item A Opção A é superior: título comunica a mensagem, uso restrito de cores direciona
  a atenção, espaço em branco organiza a leitura, e o dashboard passa no Teste do Relance.
  \item A Opção B é superior porque tem mais gráficos e mais cores, oferecendo mais informação.
  \item As duas opções são equivalentes em eficácia comunicativa.
  \item A Opção A é inferior porque tem espaço desperdiçado na tela.
  \item Títulos descritivos como ``Dashboard de Logística'' são mais profissionais.
\end{enumerate}

\textbf{Gabarito: A.}

\textbf{Justificativa:} a Opção A aplica os princípios fundamentais: título-mensagem (não
rótulo), hierarquia visual clara (1 KPI principal), uso restrito de cores (máx.\ 3 com
propósito), e espaço em branco como organizador da leitura — garantindo que o Teste do
Relance seja aprovado em 3 segundos. A Opção B viola: título rótulo (B), excesso de
visuais que fragmenta a atenção (B), cores sem propósito semântico (B) e rótulos
descritivos que não informam (B).
\end{SolvedBox}

% =============================================================================
\section{Bloco III — Tableau Interativo: Filtros, Parâmetros e Ações (S10)}
% =============================================================================

\begin{BoardBox}
\textbf{Hierarquia de filtros no Tableau — para o quadro}

O Tableau aplica filtros em uma ordem fixa. Conhecer essa ordem evita resultados
inesperados:
\begin{enumerate}
  \item \textbf{Extract filter} — filtra antes de carregar
  \item \textbf{Data source filter} — nível da fonte de dados
  \item \textbf{Context filter} — define o subconjunto base para outros filtros
  \item \textbf{Dimension filter} — filtra dimensões (categorias)
  \item \textbf{Measure / Table calculation filter} — filtra sobre métricas e cálculos
\end{enumerate}

\textbf{Armadilha clássica:} Top N sem Context Filter. O Tableau calcula o Top N sobre
o dataset inteiro, ignorando filtros de região já aplicados. Solução: marcar o filtro de
região como \textit{Context Filter} primeiro.
\end{BoardBox}

\begin{SolvedBox}
\textbf{Filtros vs Parâmetros — diferença fundamental:}

\begin{tabular}{|l|p{5cm}|p{5.5cm}|}
\hline
 & \textbf{Filtro (Quick Filter)} & \textbf{Parâmetro} \\
\hline
\textbf{O que faz} & Inclui/exclui linhas do dataset & Muda o \emph{cálculo} sem filtrar linhas \\
\textbf{Quando usar} & Selecionar região, categoria, período & Alternar métrica, definir meta, Top N \\
\textbf{Exemplos} & ``Ver apenas PE''; ``Últimos 30 dias'' &
``Qual métrica ver: Receita ou Margem?''; ``Meta de vendas: R\$ X'' \\
\hline
\end{tabular}

\textbf{Campo calculado com parâmetro de meta:}\\
\texttt{IF [Vendas] >= [Meta Vendas] THEN "Atingiu" ELSE "Abaixo" END}

O usuário ajusta \texttt{[Meta Vendas]} em tempo real — o gráfico atualiza imediatamente.
\end{SolvedBox}

\begin{BoardBox}
\textbf{Ações de Dashboard — os 4 tipos (para o quadro)}

\begin{tabular}{|l|p{5cm}|p{5.5cm}|}
\hline
\textbf{Tipo} & \textbf{O que faz} & \textbf{Exemplo de uso} \\
\hline
\textbf{Filter action} & Clicar em um visual filtra outro & Clicar em ``SP'' filtra todas as demais views \\
\textbf{Highlight action} & Destaca elementos relacionados sem filtrar & Passar o mouse em Região destaca a mesma no mapa \\
\textbf{URL action} & Abre URL com dado selecionado & Clique no produto abre a ficha no sistema ERP \\
\textbf{Navigate action} & Navega entre dashboards & Clicar em KPI leva ao painel de drill-down \\
\hline
\end{tabular}
\end{BoardBox}

% ---------------------------------------------------------------------------
\subsection{Questão Tipo Prova III — Filtros e Parâmetros (S10)}
% ---------------------------------------------------------------------------

\begin{SolvedBox}
\textbf{Questão Q-A3 — Tableau: Filtros e Parâmetros}

André é analista em uma rede de 30 postos de combustível em PE. O gerente pediu que o
dashboard permitisse:
\begin{enumerate}
  \item Ver apenas os postos de uma cidade específica (ex.: só Caruaru)
  \item Escolher o período de análise (último mês, trimestre ou ano)
  \item Alternar a métrica exibida entre Volume de Litros e Receita Total
\end{enumerate}

Qual alternativa apresenta \textbf{corretamente} o recurso do Tableau para cada necessidade?

\begin{enumerate}[label=\Alph*)]
  \item Filtro para selecionar a cidade; Parâmetro para escolher o período
  (com campo calculado que ajusta as datas); Parâmetro para alternar a métrica
  (com campo calculado \texttt{IF [Métrica]="Receita" THEN [Receita] ELSE [Litros] END})
  \item Usar apenas gráficos estáticos sem nenhuma interatividade
  \item Criar 30 dashboards separados, um para cada posto
  \item Exportar para Excel e fazer as análises manualmente
  \item Usar apenas parâmetros para todas as 3 necessidades, sem filtros
\end{enumerate}

\textbf{Gabarito: A.}

\textbf{Justificativa:} (1) filtrar por cidade é uma operação de \emph{inclusão/exclusão
de linhas} — função natural de Quick Filter; (2) o período envolve um cálculo de datas
dinâmico, que só um Parâmetro + campo calculado resolve; (3) alternar entre métricas
muda o \emph{campo calculado}, não filtra linhas — é exclusividade do Parâmetro.
A opção E está errada porque parâmetros não filtram linhas; a opção A é a única com a
divisão correta de responsabilidades entre os dois recursos.
\end{SolvedBox}

% =============================================================================
\section{Bloco IV — Dashboards Decisórios, KPIs e Pitch de 3 minutos (S11)}
% =============================================================================

\begin{BoardBox}
\textbf{As 3 zonas de um dashboard decisório — para o quadro}

\textbf{Zona 1 — Contexto:} o que está acontecendo agora?
\begin{itemize}
  \item KPIs principais com comparação à meta ou período anterior.
  \item Lido em $<$ 5 segundos. Nenhuma interpretação ainda.
\end{itemize}

\textbf{Zona 2 — Diagnóstico:} por que está acontecendo?
\begin{itemize}
  \item Drill-down dos KPIs em alerta: por produto, região, período.
  \item Permite formular hipóteses — não confirmar causas.
\end{itemize}

\textbf{Zona 3 — Recomendação:} o que fazer?
\begin{itemize}
  \item 1–2 recomendações concretas com evidências.
  \item Call to Action explícito: \emph{quem} faz \emph{o quê} até \emph{quando}.
\end{itemize}
\end{BoardBox}

\begin{SolvedBox}
\textbf{Anatomia de um bom KPI:}

Um KPI sozinho não decide nada. Um KPI bem construído tem 5 componentes:
\begin{enumerate}
  \item \textbf{Valor atual} — o número real do período
  \item \textbf{Meta (target)} — o que deveria ser
  \item \textbf{Gap} — distância entre atual e meta (valor absoluto e \%)
  \item \textbf{Tendência} — seta ou sparkline: melhorando ou piorando?
  \item \textbf{Período} — MTD, YTD, vs ano anterior
\end{enumerate}

\textbf{Anti-padrões de KPI:}
tudo sempre verde (conforto falso); métricas sem definição (``Receita'' — qual?);
granularidade excessiva (50 KPIs na tela inicial); número isolado sem contexto de comparação.
\end{SolvedBox}

\begin{BoardBox}
\textbf{Estrutura do pitch de 3 minutos — para o quadro}

3 min $\approx$ 450 palavras $\approx$ 4 blocos:

\begin{enumerate}
  \item \textbf{Situação (30s):} contexto factual que a audiência reconhece.
  \item \textbf{Complicação (45s):} o problema ou oportunidade que a situação revela.
  \item \textbf{Resolução (60s):} evidência + recomendação com número concreto.
  \item \textbf{Próximo passo (45s):} pedido específico — verbo + quem + prazo.
\end{enumerate}

\textbf{Erro mais comum:} começar com contexto genérico em vez de dado impactante.
\end{BoardBox}

% ---------------------------------------------------------------------------
\subsection{Questão Tipo Prova IV — Dashboard Decisório (S11)}
% ---------------------------------------------------------------------------

\begin{SolvedBox}
\textbf{Questão Q-A4 — Dashboards que Forçam Decisão}

Rodrigo é gerente de expansão de uma rede de academias low-cost e precisa apresentar
para o CEO se deve abrir uma unidade em Petrolina ou em Juazeiro do Norte. Ele estruturou
o dashboard em 3 seções:

\textbf{Seção 1 — Contexto:} mapa comparando as duas cidades; indicadores de mercado
(população-alvo, renda per capita, número de academias concorrentes).

\textbf{Seção 2 — Diagnóstico:} comparativo de custo de aluguel, custo de captação de
sócio, ticket médio projetado por região, taxa de retenção de unidades similares.

\textbf{Seção 3 — Recomendação:} simulação de payback nas 2 cidades; cenários otimista,
realista e pessimista; call-to-action: ``Petrolina — abertura Q3 2026, investimento de
R\$380k, ROI em 22 meses''.

O CEO tomou a decisão em 12 minutos porque todas as informações necessárias estavam
organizadas de forma lógica e a recomendação era clara e acionável.

Qual alternativa descreve \textbf{corretamente} por que a estrutura funcionou?

\begin{enumerate}[label=\Alph*)]
  \item O framework Contexto→Diagnóstico→Recomendação funciona porque primeiro estabelece
  o cenário (o que é), depois evidencia o problema/oportunidade (por que importa), e
  finaliza com ação concreta (o que fazer); os 3 cenários reduzem incerteza; a
  recomendação com prazo e ROI torna a decisão acionável.
  \item O dashboard foi eficaz apenas porque tinha gráficos visualmente bonitos.
  \item A ordem das 3 seções é irrelevante; qualquer estrutura funcionaria igualmente.
  \item Dashboards não devem incluir recomendações — apenas dados brutos.
  \item O CEO decidiu rápido porque não leu o dashboard com atenção suficiente.
\end{enumerate}

\textbf{Gabarito: A.}

\textbf{Justificativa:} a estrutura em 3 zonas é eficaz porque respeita o processo de
decisão humano: primeiro entender o cenário (Contexto), depois identificar o problema
(Diagnóstico), por fim receber uma recomendação sustentada por evidências (Recomendação
com ROI e prazo). Os cenários otimista/realista/pessimista são essenciais para gestão
de risco — uma decisão de R\$380k sem análise de cenários não seria tomada.
\end{SolvedBox}

% =============================================================================
\section{Simulado Integrado — 5 Questões no Formato das Provas}
% =============================================================================

\begin{NoteBox}
\textbf{Simulado final:} estas 5 questões cobrem todo o conteúdo da Unidade II no
formato exato das provas (Prova II Unidade, Final e II Chamada). Complete as questões
individualmente e depois discutimos os gabaritos.

\textbf{Tempo sugerido:} 30 minutos para responder + 15 minutos de correção coletiva.
\end{NoteBox}

% ---------------------------------------------------------------------------
\subsection*{Questão S1 (2,0 pts) — Storytelling com Dados}
% ---------------------------------------------------------------------------

\begin{SolvedBox}
\textbf{Cenário:}

Isabela é analista de dados em uma operadora de planos de saúde com 80 mil beneficiários
no interior de PE. Após analisar os dados de utilização do plano, ela descobriu que:
\begin{itemize}
  \item 57\% das internações hospitalares ocorrem em beneficiários com doenças crônicas
  não monitoradas (diabetes, hipertensão)
  \item Beneficiários que participam do programa de telemonitoramento têm 62\% menos
  internações do que os que não participam
  \item O custo médio de uma internação é R\$8.400; o custo mensal do telemonitoramento
  por beneficiário é R\$45
\end{itemize}

Isabela vai apresentar para o comitê de gestão de saúde. Qual alternativa representa
a \textbf{Big Idea bem formulada}?

\begin{enumerate}[label=\Alph*)]
  \item ``Devemos expandir o programa de telemonitoramento para todos os beneficiários
  com doenças crônicas, pois a análise indica redução de 62\% nas internações — economia
  estimada de R\$18,2M/ano contra um investimento de R\$3,6M no programa.''
  \item ``Os dados mostram informações relevantes sobre internações hospitalares.''
  \item ``Analisamos 80.000 registros de beneficiários no período de 2024–2025.''
  \item ``A taxa de internação é de 7,3\% ao ano.''
  \item ``O gráfico de barras mostra a distribuição de internações por tipo de doença.''
\end{enumerate}

\medskip
\textbf{Gabarito: A.}

\textbf{Justificativa detalhada:} a alternativa A cumpre os 3 critérios:
(1) \textbf{ponto de vista} — ``devemos expandir'' é uma recomendação de ação, não uma
observação passiva; (2) \textbf{tensão} — o dado evidencia o risco (57\% das internações
evitáveis) e a oportunidade quantificada (62\% menos internações); (3) \textbf{stake
financeiro} — R\$18,2M de economia vs R\$3,6M de investimento é o argumento decisivo
para o comitê de gestão. As demais alternativas são vagas~(B), descritivas sem ação~(C),
número isolado sem contexto~(D), ou mero rótulo de visual~(E).
\end{SolvedBox}

% ---------------------------------------------------------------------------
\subsection*{Questão S2 (2,0 pts) — Design de Informação}
% ---------------------------------------------------------------------------

\begin{SolvedBox}
\textbf{Cenário:}

Gustavo está criando um dashboard no Tableau para o VP de Vendas de uma rede de
franquias de alimentação. O VP tem apenas 5 minutos por dia. Gustavo está avaliando
duas abordagens de design:

\textbf{Opção 1:}
\begin{itemize}
  \item Título: ``Franchisees abaixo da meta em Abril: 8 de 24 unidades (33\%) —
  foco em lojas do interior''
  \item KPI central: porcentagem de unidades na meta (verde se $\geq$ 80\%, vermelho se $<$ 80\%)
  \item Barras horizontais com a unidade mais crítica destacada em laranja
  \item Fundo branco, espaço em branco generoso, máx.\ 3 cores com propósito semântico
\end{itemize}

\textbf{Opção 2:}
\begin{itemize}
  \item Título: ``Painel de Desempenho de Franquias''
  \item 16 gráficos de pizza em diferentes tamanhos sem ordem evidente
  \item Gradientes e sombras em todos os elementos
  \item Legenda separada com 12 entradas de cor diferente
  \item Todos os números com 4 casas decimais
\end{itemize}

Qual alternativa avalia corretamente as duas opções?

\begin{enumerate}[label=\Alph*)]
  \item A Opção 1 é superior: o título comunica a mensagem antes de o VP ler qualquer
  gráfico; o KPI com cor semântica (verde/vermelho) passa no Teste do Relance em 3
  segundos; o destaque em laranja direciona a ação; o espaço em branco organiza a
  hierarquia visual.
  \item A Opção 2 é superior: mais gráficos significa mais informação e maior credibilidade.
  \item Ambas são equivalentes — o que importa é o dado, não o design.
  \item A Opção 1 tem excesso de espaço em branco que ``desperdiça'' a tela.
  \item Títulos descritivos como ``Painel de Desempenho'' transmitem mais profissionalismo.
\end{enumerate}

\medskip
\textbf{Gabarito: A.}

\textbf{Justificativa:} a Opção 1 aplica: (a) título-mensagem com dado concreto e
relevância imediata; (b) cor semântica (verde/vermelho) como atributo pré-atencional
para status de KPI; (c) 1 cor de destaque (laranja) para direcionar ação; (d) Data-ink
ratio maximizado (sem chartjunk). A Opção 2 viola: o gráfico de pizza com 12 fatias
impede comparação precisa (ângulos vs comprimentos); gradientes são chartjunk (Tufte);
legenda separada com 12 entradas exige desvio do olhar; 4 casas decimais são precisão
falsa para dados de negócio.
\end{SolvedBox}

% ---------------------------------------------------------------------------
\subsection*{Questão S3 (2,0 pts) — Tableau: Filtros e Parâmetros}
% ---------------------------------------------------------------------------

\begin{SolvedBox}
\textbf{Cenário:}

Fernanda é analista de BI em uma rede de clínicas veterinárias com 25 unidades em
Pernambuco e Alagoas. A gerência pediu que o dashboard permitisse:
\begin{enumerate}
  \item Visualizar apenas as clínicas de um estado específico (PE ou AL)
  \item Definir dinamicamente a meta de consultas mensais (o número muda todo trimestre)
  \item Escolher entre visualizar Receita Total ou Margem Líquida
\end{enumerate}

Qual alternativa apresenta \textbf{corretamente} o recurso do Tableau para cada necessidade?

\begin{enumerate}[label=\Alph*)]
  \item Filtro por estado; Parâmetro para a meta (campo calculado: \\
  \texttt{IF [Consultas] >= [Meta Consultas] THEN "OK" ELSE "Alerta" END}); \\
  Parâmetro de métrica (campo calculado que alterna entre Receita e Margem)
  \item Criar 25 dashboards separados, um por clínica
  \item Usar apenas filtros para as 3 necessidades, sem parâmetros
  \item Exportar os dados para Excel e calcular manualmente
  \item Publicar três versões diferentes do dashboard, uma para cada necessidade
\end{enumerate}

\medskip
\textbf{Gabarito: A.}

\textbf{Justificativa:} (1) filtrar por estado é \emph{inclusão/exclusão de linhas} —
responsabilidade do Quick Filter; (2) a meta muda o \emph{valor de referência do cálculo},
não filtra dados — é responsabilidade do Parâmetro + campo calculado; (3) alternar
entre Receita e Margem muda \emph{qual campo calculado} é exibido, operação que só o
Parâmetro realiza. O filtro não consegue fazer (2) nem (3) porque não altera a lógica
do cálculo — apenas inclui ou exclui registros.
\end{SolvedBox}

% ---------------------------------------------------------------------------
\subsection*{Questão S4 (2,0 pts) — Dashboards Decisórios}
% ---------------------------------------------------------------------------

\begin{SolvedBox}
\textbf{Cenário:}

Mariana é diretora de TI em uma empresa de logística. Ela precisa convencer o board
a aprovar a migração do sistema de rastreamento para a nuvem. Montou o dashboard em
3 seções:

\textbf{Seção 1 — Contexto:} tempo médio de inatividade atual (14h/mês), custo
operacional do sistema legado (R\$480k/ano), satisfação dos operadores (NPS interno: 31).

\textbf{Seção 2 — Diagnóstico:} análise das causas de inatividade (78\% por falhas
de hardware com mais de 7 anos); benchmarking com concorrentes; impacto de cada hora
de inatividade (R\$38k em SLAs não cumpridas).

\textbf{Seção 3 — Recomendação:} projeção de custos na nuvem vs manutenção do legado
em 3 anos; ROI de 240\% em 18 meses; call-to-action: ``Aprovação de R\$280k até
31/05/2026 para início da migração em junho''.

O board aprovou o projeto na reunião. Qual alternativa explica corretamente o porquê?

\begin{enumerate}[label=\Alph*)]
  \item O framework C→D→R é eficaz: o Contexto estabelece a urgência com dados reais;
  o Diagnóstico elimina dúvidas sobre a causa raiz; a Recomendação reduz incerteza
  com ROI e prazo; o Call to Action com data e valor exato torna a aprovação acionável.
  \item O board aprovou porque os gráficos eram bonitos e coloridos.
  \item A estrutura das 3 seções é irrelevante; qualquer ordem funcionaria igual.
  \item Dashboards não deveriam incluir recomendações — o board deveria chegar às
  próprias conclusões com dados brutos.
  \item O projeto foi aprovado porque o orçamento pedido era pequeno.
\end{enumerate}

\medskip
\textbf{Gabarito: A.}

\textbf{Justificativa:} o framework Contexto→Diagnóstico→Recomendação funciona porque
replica o raciocínio decisório humano. O Contexto com métricas quantificadas (14h/mês,
NPS 31) cria urgência sem parecer alarmismo. O Diagnóstico com causa raiz (hardware
com 7+ anos, 78\% das falhas) fecha a possibilidade de ``talvez seja outra coisa''.
A Recomendação com ROI de 240\% em 18 meses e call-to-action com data transforma uma
apresentação informativa em um dashboard decisório que \emph{força} a decisão.
\end{SolvedBox}

% ---------------------------------------------------------------------------
\subsection*{Questão S5 (2,0 pts) — Pitch de 3 minutos}
% ---------------------------------------------------------------------------

\begin{SolvedBox}
\textbf{Cenário:}

Carlos apresentou sua análise sobre endividamento de microempresários na Feira de
Soluções Financeiras. Ele tinha 3 minutos para convencer os jurados de que sua análise
gera insights acionáveis.

\textbf{Minuto 1 — Situação:}
``A cada ano, 35 mil microempresários no Nordeste entram em processo de insolvência.
Isso representa R\$420 milhões em crédito perdido e 35 mil famílias sem renda.
Nossa análise identificou os 3 principais preditores de falência.''

\textbf{Minuto 2 — Complicação + Dados:}
\begin{itemize}
  \item Visual 1: mapa de calor das cidades com maior índice de insolvência
  \item Visual 2: correlação entre inadimplência $>$ 60 dias e probabilidade de falência
  \item Insight: ``microempresários com histórico de inadimplência $>$ 60 dias têm 5$\times$
  mais chance de insolvência nos 12 meses seguintes''
\end{itemize}

\textbf{Minuto 3 — Resolução + Próximo passo:}
``Um sistema de alerta precoce pode reduzir insolvências em 28\%, beneficiando
9.800 famílias. Próximo passo: piloto com 200 empresas no segundo semestre.''

Carlos termina dentro do tempo. Os jurados entendem o problema, os insights e a
proposta. Qual alternativa descreve \textbf{corretamente} os elementos eficazes do pitch?

\begin{enumerate}[label=\Alph*)]
  \item O pitch foi eficaz: gancho com dado impactante capturou atenção; apenas 2 visuais
  evitaram sobrecarga; o insight foi traduzido em ação concreta com métrica de impacto;
  o próximo passo tem responsável e prazo (estrutura Situação→Complicação→Resolução→Ação).
  \item O pitch foi eficaz apenas porque durou exatamente 3 minutos, independente do conteúdo.
  \item Carlos deveria ter mostrado todas as 40 visualizações criadas durante a análise.
  \item Pitches de dados não devem ter gancho emocional — devem começar com tabelas.
  \item A recomendação final é desnecessária; os jurados deveriam decidir sozinhos.
\end{enumerate}

\medskip
\textbf{Gabarito: A.}

\textbf{Justificativa:} o pitch de Carlos aplica a estrutura SCR (Situação, Complicação,
Resolução) de forma precisa. A Situação começa com dado quantificado que torna o problema
concreto e relevante — não com contexto genérico. A Complicação usa apenas 2 visuais
estratégicos (respeita a capacidade atencional da audiência) e entrega o insight com
magnitude (5$\times$). A Resolução tem métrica de impacto (28\%) e o próximo passo tem
verbo de ação (``piloto''), quantidade (200 empresas) e prazo (segundo semestre) —
os 3 elementos de um call-to-action eficaz.
\end{SolvedBox}

% =============================================================================
\section{Mapa de Conteúdo por Prova — Resumo Rápido}
% =============================================================================

\begin{BoardBox}
\textbf{O que cai em cada prova:}

\begin{tabular}{|l|p{4cm}|p{7cm}|}
\hline
\textbf{Prova} & \textbf{Foco} & \textbf{Tópicos da Unidade II} \\
\hline
\textbf{Prova II Unidade} & 5 questões / 2 pts cada & Big Idea; Design de Dashboard; Filtros e Parâmetros;
Dashboard Decisório; Pitch 3 min \\
\hline
\textbf{Prova Final} & 10 questões / 1 pt cada & Unidade I (Q1–Q7) + Big Idea (Q8) +
Design de Dashboard (Q9) + Filtros (Q10) \\
\hline
\textbf{Prova II Chamada} & 10 questões / 1 pt cada & Unidade I + II completas;
Big Idea; Dashboard; Filtros/Parâmetros \\
\hline
\end{tabular}

\textbf{Para a Prova Final e II Chamada:} revisar também o conteúdo da Unidade I
(dicionário de dados, qualidade, média/mediana/moda, desvio-padrão, boxplot/IQR,
dados ausentes, correlação ≠ causalidade).
\end{BoardBox}

\begin{NoteBox}
\textbf{Checklist de preparação para a prova:}
\begin{itemize}
  \item[$\square$] Sei escrever uma Big Idea com ponto de vista + tensão + stake
  \item[$\square$] Sei aplicar o Teste do Relance e identificar suas 4 causas de falha
  \item[$\square$] Sei distinguir filtro de parâmetro com exemplo concreto
  \item[$\square$] Sei descrever as 3 zonas de um dashboard decisório
  \item[$\square$] Sei estruturar um pitch de 3 minutos (Situação → Complicação → Resolução → Ação)
  \item[$\square$] Sei calcular IQR e limites de outlier (para Prova Final e II Chamada)
  \item[$\square$] Sei distinguir correlação de causalidade com exemplo
\end{itemize}
\end{NoteBox}

% =============================================================================
\section{Orientações para a Aula de 18/05 — Apresentações em Grupo}
% =============================================================================

\begin{BoardBox}
\textbf{Semana 18/05 — Apresentação em Grupo: Dashboard Storytelling}

\begin{itemize}
  \item \textbf{Dinâmica:} cada grupo terá um dataset do Kaggle, uma dor de negócio e
  uma pergunta norteadora.
  \item \textbf{Entrega:} dashboard funcional no Tableau Public + pitch de 3 minutos
  na estrutura SCR.
  \item \textbf{Estrutura do dashboard:} Pergunta → Big Number (esq.) + Insight IA (dir.)
  → Gráficos que respondem → Próxima pergunta → … → Tabela Analítica Final.
  \item \textbf{Insight IA:} cada grupo deve gerar um insight textual usando ChatGPT/Gemini
  a partir dos dados, e colocá-lo ao lado do Big Number como um painel de texto no
  dashboard.
  \item \textbf{Critério principal:} o dashboard \emph{força uma decisão} ou apenas
  exibe dados?
\end{itemize}

\textbf{Os grupos e datasets serão anunciados na aula de 18/05. Protótipos já
disponíveis para consulta (ver \texttt{semana12\_grupos\_tableau.tex}).}
\end{BoardBox}
